% ══════════════════════════════════════════════════════════════════
% CHAPITRE 3 — IMPLÉMENTATION ET TESTS
% Plan : Intro · 3.1 → 3.8 · Conclusion
% ══════════════════════════════════════════════════════════════════

\chapter{Implémentation et Tests}
\label{chap:implementation}

\section*{Introduction}
\addcontentsline{toc}{section}{Introduction}

Ce chapitre décrit la réalisation concrète du système conçu au
chapitre précédent. Il présente successivement l'environnement de
développement, la collecte et le prétraitement des données,
l'implémentation des trois modules d'Intelligence Artificielle,
du module Historique et de l'application web qui les intègre,
puis définit le protocole de tests appliqué à l'ensemble.

% ══════════════════════════════════════════════════════════════════
\section{Environnement de développement}
\label{sec:environnement}
% ══════════════════════════════════════════════════════════════════

Deux environnements complémentaires ont été utilisés : la
plateforme cloud Kaggle pour l'entraînement des modèles
(accélération GPU), et un poste local sous Ubuntu pour le
développement et les tests de l'application. Le
tableau~\ref{tab:environnements} en résume les caractéristiques.

\begin{table}[H]
	\centering
	\caption{Environnements de développement utilisés}
	\label{tab:environnements}
	\small
	\begin{tabularx}{\textwidth}{lXX}
		\toprule
		\textbf{Composant} & \textbf{Entraînement (Kaggle)} &
		\textbf{Développement (local)} \\
		\midrule
		Système       & Linux (cloud)         & Ubuntu 24.04 LTS \\
		Processeur    & 2 GPU Tesla T4        & CPU Intel (HP EliteBook) \\
		Python        & 3.12                  & 3.12 \\
		TensorFlow    & 2.20 (entraînement)   & TFLite Runtime 2.14 \\
		Scikit-learn  & 1.3+                  & 1.3+ \\
		XGBoost       & 2.0+                  & 2.0+ \\
		Framework web & ---                   & Streamlit 1.35+ \\
		SGBD          & ---                   & SQLite 3 (natif Python) \\
		Éditeur       & Kaggle Notebook       & VS Code \\
		Versionnage   & ---                   & Git / GitHub \\
		\bottomrule
	\end{tabularx}
\end{table}

L'application est organisée selon la structure multi-pages de
Streamlit :

\begin{lstlisting}[language=bash,
	caption={Structure des fichiers de l'application},
	label={lst:structure}]
	PFE_Agriculture_IA/
	|-- app.py                  # Page d'accueil
	|-- database.py             # Module Historique (SQLite)
	|-- meteo_api.py            # Client API Open-Meteo
	|-- styles.py               # Feuille de style globale
	|-- requirements.txt        # Dependances Python
	|-- models/
	|   |-- model_irrigation.pkl       # Random Forest
	|   |-- scaler_irrigation.pkl
	|   |-- model_fertilisation.pkl    # XGBoost
	|   |-- scaler_fertilisation.pkl
	|   |-- model_maladie.tflite       # CNN MobileNetV2
	|   `-- config_maladies.json       # Liste des classes
	`-- pages/
	|-- 1_Irrigation.py
	|-- 2_Fertilisation.py
	|-- 3_Detection_Maladies.py
	`-- 4_Historique.py
\end{lstlisting}

% ══════════════════════════════════════════════════════════════════
\section{Collecte et prétraitement des données}
\label{sec:collecte}
% ══════════════════════════════════════════════════════════════════

\subsection{Données d'irrigation --- NASA POWER}

Les données climatiques ont été extraites de l'API NASA POWER
\cite{stackhouse2019nasa} pour les cinq villes cibles
(N'Djamena, Abéché, Sarh, Moundou, Bongor) sur la période
2015--2023, soit \textbf{16~435 observations journalières}.
Quinze variables explicatives ont été retenues : température
moyenne et maximale, cumuls de pluie sur 7 et 14 jours, humidité
de l'air et du sol, ETP, vent, rayonnement, ainsi que des
variables agronomiques encodées (culture, stade phénologique,
saison, type de sol, pH, zone climatique). La variable cible
binaire (irriguer : OUI/NON) a été construite à partir de règles
agronomiques croisant bilan hydrique et stade de la culture.

Le prétraitement a comporté quatre étapes : encodage des
variables catégorielles (LabelEncoder), normalisation des
variables numériques (StandardScaler), traitement des valeurs
manquantes par interpolation, et partitionnement stratifié
en ensembles d'entraînement (70\,\%), de validation (15\,\%) et
de test (15\,\%).

\subsection{Données de fertilisation}

Le jeu de données de fertilisation combine 99 observations
réelles issues d'un dataset public Kaggle et 1~500 observations
générées par augmentation cohérente (bruit gaussien contrôlé
autour des signatures NPK de chaque engrais), soit
\textbf{1~599 observations} couvrant cinq classes d'engrais :
Urée, DAP, NPK 15-15-15, NPK 23-10-5 et KCl.

\subsection{Dataset local d'images --- maladies}

Un dataset de \textbf{287 images} de feuilles de cultures
tchadiennes a été constitué, réparti en \textbf{10 classes}
(5 cultures $\times$ 2 états : saine/malade), comme détaillé en
annexe~A. Les images, prises en conditions réelles de terrain,
ont été redimensionnées en $224 \times 224$ pixels et normalisées
dans $[0,1]$. Pour compenser la petite taille du dataset, une
augmentation de données a été appliquée à l'entraînement :
rotations ($\pm 30°$), symétries horizontale et verticale, zoom
($\pm 20\,\%$) et variations de luminosité ($0{,}65$ à $1{,}35$).

% ══════════════════════════════════════════════════════════════════
\section{Module IA --- Irrigation}
\label{sec:module_irrigation}
% ══════════════════════════════════════════════════════════════════

\subsection{Choix de l'algorithme Random Forest}

Le Random Forest \cite{breiman2001random} agrège les prédictions
de $T$ arbres de décision entraînés sur des échantillons
\textit{bootstrap} ; la classe finale résulte du vote
majoritaire :

\begin{equation}
	\hat{y} = \arg\max_{c \in \{0,1\}} \sum_{t=1}^{T}
	\mathbf{1}\left[ h_t(\mathbf{x}) = c \right]
	\label{eq:rf}
\end{equation}

où $h_t(\mathbf{x})$ désigne la prédiction du $t$-ième arbre
pour l'observation $\mathbf{x}$. Cet algorithme a été retenu
pour sa robustesse au bruit, sa résistance au surapprentissage
et sa capacité à mesurer l'importance des variables.

\subsection{Entraînement et hyperparamètres}

Les hyperparamètres ont été optimisés par validation croisée à 5
plis. La configuration retenue est : 100 arbres
(\texttt{n\_estimators}), profondeur non limitée
(\texttt{max\_depth=None}), pondération équilibrée des classes
(\texttt{class\_weight=balanced}) et graine fixée
(\texttt{random\_state=42}) pour la reproductibilité.

\subsection{Intégration dans l'application}

Le modèle et son normaliseur, sérialisés avec \texttt{joblib},
sont chargés une seule fois grâce au cache de Streamlit :

\begin{lstlisting}[language=Python,
	caption={Chargement et prédiction du module Irrigation},
	label={lst:pred_irr}]
	import joblib, numpy as np
	
	model  = joblib.load("models/model_irrigation.pkl")
	scaler = joblib.load("models/scaler_irrigation.pkl")
	
	def predire_irrigation(features):
	X  = scaler.transform(np.array([features]))
	y  = model.predict(X)[0]
	p  = model.predict_proba(X)[0]
	return {"decision": "OUI" if y == 1 else "NON",
		"confiance": round(p[y] * 100, 1)}
\end{lstlisting}

% ══════════════════════════════════════════════════════════════════
\section{Module IA --- Fertilisation}
\label{sec:module_fertilisation}
% ══════════════════════════════════════════════════════════════════

\subsection{Choix de l'algorithme XGBoost}

XGBoost \cite{chen2016xgboost} construit séquentiellement des
arbres correcteurs en minimisant une fonction objectif
régularisée :

\begin{equation}
	\mathcal{L}(\phi) = \sum_{i=1}^{n} l\left(y_i, \hat{y}_i\right)
	+ \sum_{k=1}^{K} \Omega\left(f_k\right)
	\label{eq:xgb}
\end{equation}

où $l$ est la perte (entropie croisée multi-classes) et
$\Omega(f_k)$ pénalise la complexité de l'arbre $k$, limitant le
surapprentissage --- un atout essentiel sur notre jeu de données
de taille modeste.

\subsection{Entrées et sorties du module}

Le modèle prend en entrée onze variables (N, P, K, pH,
température, humidités de l'air et du sol, pluie, culture, type
de sol, zone) et prédit l'un des cinq engrais. La recommandation
affichée est enrichie d'une dose indicative (kg/ha) et de la
période d'application, issues d'une table de correspondance
agronomique.

% ══════════════════════════════════════════════════════════════════
\section{Module IA --- Détection des maladies}
\label{sec:module_maladies}
% ══════════════════════════════════════════════════════════════════

\subsection{Architecture CNN MobileNetV2 et transfert d'apprentissage}

Le réseau retenu est MobileNetV2 \cite{sandler2018mobilenetv2}
pré-entraîné sur ImageNet, surmonté d'une tête de classification
adaptée : GlobalAveragePooling, couche dense de 256 neurones
(ReLU) avec BatchNormalization et Dropout (0,5), couche dense de
128 neurones (ReLU) avec BatchNormalization et Dropout (0,4), et
couche de sortie Softmax à 10 classes.

L'entraînement, mené sur Kaggle (2 GPU T4), a suivi une stratégie
en deux phases : (i) base gelée, seule la tête est entraînée
(taux d'apprentissage $10^{-3}$, arrêt précoce avec patience de
10 époques) ; (ii) \textit{fine-tuning} des 30 dernières couches
de la base (taux réduit à $10^{-4}$).

\subsection{Conversion au format TensorFlow Lite}

Pour garantir un déploiement léger et compatible avec
l'environnement cloud, le modèle Keras a été converti au format
TFLite \cite{tensorflow2023} avec quantification par défaut,
ramenant sa taille à moins de 10~Mo. L'inférence dans
l'application s'effectue via la bibliothèque légère
\texttt{tflite-runtime}, sans dépendance à TensorFlow complet.

\subsection{Analyse NDVI : satellite et caméra smartphone}

En complément du diagnostic par photographie, le module propose
une évaluation par indice de végétation selon deux sources :

\begin{itemize}[leftmargin=1.2cm]
	\item \textbf{Satellite Sentinel-2} \cite{sentinel2esa} :
	l'utilisateur relève la valeur NDVI de sa parcelle sur
	la plateforme gratuite EO Browser et la saisit dans
	l'application, qui l'interprète selon des seuils
	spécifiques à chaque culture ;
	\item \textbf{Caméra smartphone} : à défaut de capteur proche
	infrarouge, l'indice VARI est calculé sur les canaux
	visibles de la photographie :
\end{itemize}

\begin{equation}
	\text{VARI} = \frac{G - R}{G + R - B}
	\label{eq:vari}
\end{equation}

puis converti en estimation du NDVI par la relation empirique :

\begin{equation}
	\text{NDVI}_{\text{estimé}} = 0{,}18 + 0{,}73 \times \text{VARI}
	\label{eq:vari_ndvi}
\end{equation}

Le diagnostic final croise l'indice obtenu avec les conditions
climatiques (température, humidité, pluie) pour classer la
parcelle en trois états : saine, en stress ou malade, assortis
d'un conseil de traitement.

% ══════════════════════════════════════════════════════════════════
\section{Module Historique --- Persistance SQLite et export CSV}
\label{sec:module_historique}
% ══════════════════════════════════════════════════════════════════

\subsection{Implémentation de la couche d'accès aux données}

Le module \texttt{database.py} centralise toutes les opérations
sur la base \texttt{historique.db}, conformément au modèle défini
à la section~\ref{sec:bdd_historique}. À l'initialisation, les
trois tables sont créées si elles n'existent pas
(\texttt{CREATE TABLE IF NOT EXISTS}). Chaque module d'analyse
appelle ensuite sa fonction de sauvegarde dédiée :

\begin{lstlisting}[language=Python,
	caption={Sauvegarde d'une analyse d'irrigation dans SQLite},
	label={lst:sqlite}]
	import sqlite3
	from datetime import datetime
	
	def sauvegarder_irrigation(ville, culture, resultat,
	dose_mm, confiance):
	conn = sqlite3.connect("historique.db")
	conn.execute("""
	INSERT INTO irrigation
	(date_analyse, ville, culture,
	resultat, dose_mm, confiance)
	VALUES (?, ?, ?, ?, ?, ?)""",
	(datetime.now().strftime("%Y-%m-%d %H:%M"),
	ville, culture, resultat, dose_mm, confiance))
	conn.commit()
	conn.close()
\end{lstlisting}

L'utilisation de requêtes paramétrées (\texttt{?}) protège la
base contre les injections SQL.

\subsection{Consultation, statistiques et export}

La page Historique agrège les trois tables et offre :

\begin{itemize}[leftmargin=1.2cm]
	\item une \textbf{vue chronologique unifiée} des analyses
	(module, culture, ville, résultat, confiance) ;
	\item des \textbf{statistiques globales} : nombre d'analyses
	par module, culture la plus analysée, zone à risque
	(ville comptant le plus de diagnostics \og malade \fg),
	taux d'irrigation recommandée et taux de maladies
	détectées ;
	\item l'\textbf{export CSV} de chaque table ou de l'historique
	complet, généré en mémoire et proposé au téléchargement ;
	\item la \textbf{purge sélective} de l'historique, par module
	ou en totalité.
\end{itemize}

% ══════════════════════════════════════════════════════════════════
\section{Développement de l'application web}
\label{sec:appli_web}
% ══════════════════════════════════════════════════════════════════

\subsection{Framework Streamlit}

L'interface a été développée avec Streamlit \cite{streamlit2024},
framework Python permettant de construire des applications web de
données sans développement front-end distinct. L'application
compte cinq pages : l'accueil (métriques des modèles et
statistiques de l'historique) et les quatre modules. Les modèles
sont chargés une seule fois par session grâce au décorateur
\texttt{@st.cache\_resource}, garantissant des temps de réponse
inférieurs à la contrainte des 3 secondes.

\subsection{Intégration de la météo temps réel}

Le module \texttt{meteo\_api.py} interroge l'API Open-Meteo
\cite{zippenfenig2023openmeteo} pour les cinq villes couvertes :
conditions courantes (température, humidité, vent), cumuls de
pluie sur 7 et 14 jours, ETP et humidité du sol. Les réponses
sont mises en cache 30 minutes (\texttt{@st.cache\_data}) afin de
limiter les appels réseau. Ces valeurs pré-remplissent
automatiquement les formulaires d'analyse, réduisant la saisie
manuelle de l'agriculteur.

\subsection{Interface et expérience utilisateur}

Une feuille de style commune (\texttt{styles.py}) assure la
cohérence visuelle : thème aux couleurs de la végétation,
cartes de résultats colorées selon la décision (bleu :
irrigation recommandée ; vert : parcelle saine ; orange :
stress ; rouge : maladie), et affichage systématique de la
confiance de la prédiction pour une décision éclairée.

% ══════════════════════════════════════════════════════════════════
\section{Protocole de tests}
\label{sec:protocole}
% ══════════════════════════════════════════════════════════════════

\subsection{Métriques d'évaluation des modèles}

Les modèles de classification sont évalués à partir de la matrice
de confusion (vrais/faux positifs $TP$, $FP$ ; vrais/faux
négatifs $TN$, $FN$) au moyen des métriques usuelles :

\begin{equation}
	\text{Accuracy} = \frac{TP + TN}{TP + TN + FP + FN}
	\label{eq:accuracy}
\end{equation}

\begin{equation}
	\text{Précision} = \frac{TP}{TP + FP}
	\qquad
	\text{Rappel} = \frac{TP}{TP + FN}
	\label{eq:prec_rec}
\end{equation}

\begin{equation}
	\text{F1} = 2 \times
	\frac{\text{Précision} \times \text{Rappel}}
	{\text{Précision} + \text{Rappel}}
	\label{eq:f1}
\end{equation}

S'y ajoutent l'\textbf{AUC-ROC} (aire sous la courbe ROC,
mesurant le pouvoir discriminant indépendamment du seuil de
décision) et le \textbf{score de validation croisée} à $k=5$
plis, qui apprécie la stabilité du modèle.

\subsection{Démarche de validation des modèles}

Chaque jeu de données a été partitionné de façon stratifiée en
entraînement (70\,\%), validation (15\,\%) et test (15\,\%).
L'optimisation des hyperparamètres s'est appuyée exclusivement
sur les ensembles d'entraînement et de validation ; les
performances rapportées au chapitre suivant sont mesurées sur
l'ensemble de test, jamais vu pendant l'entraînement.

\subsection{Plan de tests fonctionnels de l'application}

Au-delà des modèles, l'application elle-même a été soumise à un
plan de tests fonctionnels couvrant : le pré-remplissage météo
(BF1), le déroulement complet d'une analyse pour chacun des
trois modules d'IA (BF2 à BF5), et le module Historique ---
sauvegarde automatique, consultation unifiée, exactitude des
statistiques, export CSV et purge sélective (BF6 à BF8). Ces
tests ont été exécutés en local puis rejoués sur l'environnement
de production ; leurs résultats sont présentés au
chapitre~\ref{chap:resultats}.

\section*{Conclusion}
\addcontentsline{toc}{section}{Conclusion}

Ce chapitre a détaillé l'implémentation complète du système :
préparation de trois jeux de données, entraînement des modèles
Random Forest, XGBoost et CNN MobileNetV2 (converti en TFLite),
réalisation du module Historique fondé sur SQLite avec export
CSV, et intégration de l'ensemble dans une application web
Streamlit alimentée par la météo temps réel. Un protocole de
tests rigoureux --- métriques, partitionnement stratifié et plan
de tests fonctionnels --- a été défini. Le chapitre suivant en
présente les résultats et les discute.